\documentclass[11pt, a4paper]{article}

\usepackage[margin=2.5cm]{geometry}
\usepackage{graphicx}
\usepackage{booktabs}
\usepackage{hyperref}
\usepackage{xcolor}
\usepackage{amsmath}
\usepackage{amssymb}
\usepackage{caption}
\usepackage{subcaption}
\usepackage{float}
\usepackage{csvsimple}
\usepackage{datetime2}
\usepackage{enumitem}
\usepackage[title]{appendix}
\usepackage{multirow}

\graphicspath{{../../figures/}}

\hypersetup{
    colorlinks=true,
    linkcolor=blue!60!black,
    citecolor=green!50!black,
    urlcolor=blue!70!black,
}

\title{EEG Tokenizer Comparison Journal\\[0.3em]
\large Foundry — AJILE12 Experiments}
\author{Foundry Team}
\date{Last updated: \today}

\begin{document}
\maketitle

\begin{abstract}
    This document tracks ongoing experiments comparing different EEG tokenizer
    architectures within the Foundry framework. The goal is to determine
    whether the current baseline tokenizer---a simple per-channel,
    per-timepoint linear projection---can be improved upon, particularly
    for neural foundation models that must handle variable sampling rates
    and variable channel counts.

    Two orthogonal design axes are explored: the \emph{channel strategy}
    (how multi-channel signals are handled before temporal processing) and
    the \emph{temporal embedding} (how the time-domain signal is mapped to
    tokens). Experiments use the AJILE12 intracranial ECoG dataset
    (12~subjects) with the POYO-EEG Perceiver-IO backbone, varying only the
    tokenizer. Two downstream tasks are evaluated: 5-class activity behavior
    classification (AUROC) and pose estimation regression ($R^2$).

    This is a living document. Conclusions are tentative and reflect the
    current state of an ongoing investigation.
\end{abstract}

\tableofcontents
\newpage

%% ===================================================================
\section{Introduction}
\label{sec:introduction}
%% ===================================================================

A central challenge in building neural foundation models for
electroencephalography (EEG) and electrocorticography (ECoG) is handling
the diversity of recording setups encountered in practice. Datasets differ
in sampling rate (commonly 128--2048\,Hz), channel count (from a handful
of scalp electrodes to hundreds of intracranial contacts), and electrode
placement. An ideal tokenizer should ingest all of these without
requiring resampling to a canonical frequency or a fixed channel layout.

\subsection{The Baseline: Per-Channel Per-Timepoint Linear}

Our previous approach treats each channel independently and applies a
simple linear projection at each timepoint, producing one token per
channel per timepoint. While this satisfies the variable-input
requirements---it works for any sampling rate and any number of
channels---it has significant drawbacks:

\begin{itemize}[nosep]
    \item \textbf{Compute scaling.} Token count grows as
          $C \times T$, where $C$ is the number of channels and
          $T$ is the number of timepoints. For high-density recordings
          (e.g., 128 channels at 1000\,Hz), a single 1-second window
          produces 128{,}000 tokens---an enormous burden on the Perceiver
          backbone's attention mechanism.
    \item \textbf{Ignoring temporal structure.} Each timepoint is projected
          independently, discarding the rich temporal correlation structure
          that is fundamental to neural signals.
    \item \textbf{Ignoring spatial correlation.} Each channel is tokenized
          independently, even though neighboring electrodes in EEG/ECoG
          often carry highly correlated signals.
\end{itemize}

In practice, we had trained large models with this baseline and did not
observe the expected scaling benefits, motivating a systematic exploration
of alternatives.

\subsection{Motivation}

The experiments documented here explore two orthogonal axes of improvement:

\begin{enumerate}[nosep]
    \item \textbf{Temporal embedding:} Can we replace the per-timepoint
          linear projection with something that captures temporal structure
          while remaining sampling-rate agnostic? Two candidates are
          investigated:
          \begin{itemize}[nosep]
              \item A \emph{learnable Continuous Wavelet Transform} (CWT),
                    which decomposes the signal into frequency bands using
                    differentiable Morlet wavelets defined in physical
                    units (Hz, seconds)---making it inherently
                    sampling-rate invariant.
              \item A \emph{ResampleCNN}, which first resamples to a
                    fixed token rate via differentiable interpolation
                    (with anti-aliasing), then applies a 1D convolutional
                    stack.
          \end{itemize}
    \item \textbf{Channel strategy:} Can we reduce the per-channel token
          explosion by mixing channels \emph{before} temporal processing?
          The hypothesis is that a learned spatial projection can capture
          most of the relevant variance while dramatically reducing token
          count.
\end{enumerate}

The goal is not to prove a specific claim, but to verify whether the
baseline is suboptimal and to characterize which alternatives are most
promising for future development.

\subsection{Token Count and Compute}

Figure~\ref{fig:token-scaling} illustrates the compute implications of
different tokenizer designs. The per-timepoint linear approach produces
tokens proportional to both sampling rate and channel count, while the
CWT/CNN approaches decouple token count from sampling rate (via a
configurable target token rate), and spatial session further decouples
it from channel count.

\begin{figure}[H]
    \centering
    \includegraphics[width=0.85\textwidth]{token_count_scaling.pdf}
    \caption{Token count scaling for a 1\,s window with 64 channels.
    Per-timepoint linear scales linearly with sampling rate and channel
    count. CWT/CNN-based tokenizers fix the temporal token rate (here
    100\,Hz), and spatial session further reduces to a single stream
    of tokens independent of channel count.}
    \label{fig:token-scaling}
\end{figure}


%% ===================================================================
\section{Experimental Setup}
\label{sec:setup}
%% ===================================================================

\subsection{Dataset}

All experiments use the AJILE12 dataset (Peterson \& Brunton, 2022),
which contains intracranial ECoG recordings from 12 subjects performing
naturalistic behaviors. The dataset is well-suited for this comparison
because it has a moderate number of channels per subject, a consistent
sampling rate, and two well-defined downstream tasks (behavior
classification and pose estimation).

\subsection{Backbone}

The backbone is held constant across all experiments: a POYO-EEG model
based on the Perceiver-IO architecture with session embeddings and
rotary time embeddings. The \emph{only} component that varies is the
tokenizer module.

\subsection{Common Configuration}

\begin{table}[H]
    \centering
    \begin{tabular}{ll}
        \toprule
        \textbf{Parameter} & \textbf{Value} \\
        \midrule
        Dataset            & AJILE12 (Peterson \& Brunton, 2022), 12 ECoG subjects \\
        Split              & Intrasession (train/val from same sessions) \\
        Sequence length    & 1.0\,s \\
        Backbone           & POYOEEGModel (depth=4, heads=8, dim\_head=128) \\
        Default embed\_dim & 256 \\
        Precision          & bf16-mixed \\
        Early stopping     & patience=20 on validation metric \\
        Folds              & 2 (fold 0 and fold 1) \\
        Seed               & 42 \\
        \bottomrule
    \end{tabular}
    \caption{Shared configuration across all tokenizer experiments.}
    \label{tab:common-config}
\end{table}


%% ===================================================================
\section{Tokenizer Design Space}
\label{sec:design-space}
%% ===================================================================

The tokenizer consists of two orthogonal components: a \emph{channel
strategy} that determines how multi-channel input is handled, and a
\emph{temporal embedding} that maps the (possibly projected) time-domain
signal into token representations.

\subsection{Channel Strategies}

\subsubsection{Per-Channel}
Each channel is tokenized independently. The temporal embedding receives
a single-channel signal and produces tokens; a channel embedding
(identifying which electrode the signal came from) is fused with the
temporal tokens. This preserves per-electrode identity but produces
$C \times N_\text{time}$ tokens per window.

Two fusion variants exist:
\begin{itemize}[nosep]
    \item \textbf{Concat:} Channel embedding (dimension $d_\text{ch}$) is
          concatenated with temporal features (dimension
          $d_\text{embed} - d_\text{ch}$).
    \item \textbf{Add:} Channel embedding and temporal features both have
          dimension $d_\text{embed}$ and are summed.
\end{itemize}

\subsubsection{Spatial Session}
A learned per-session linear projection mixes all channels into a fixed
number of latent ``source'' signals before any temporal processing. Each
recording session has its own projection matrix
($C_\text{session} \to N_\text{sources}$), allowing the model to
adapt to different electrode configurations. This produces only
$N_\text{time}$ tokens per window, independent of channel count.

\textbf{Hypothesis:} ECoG channels are spatially correlated, so a linear
combination should capture most of the relevant variance while
dramatically reducing token count. The potential downside is loss of
fine-grained per-electrode information.

Variants include an optional shared ``common layer'' applied after the
per-session projection, and an MLP (2-layer) projector instead of a
single linear layer.

\subsection{Temporal Embeddings}

\subsubsection{Learnable CWT}
A differentiable Continuous Wavelet Transform using complex Morlet
wavelets at \emph{learnable} center frequencies and cycle counts.
Produces a time--frequency representation (magnitude and phase) that is
resampled to a fixed number of time tokens via \texttt{grid\_sample}.

Key properties:
\begin{itemize}[nosep]
    \item \textbf{Sampling-rate invariant} by construction: wavelets are
          defined in Hz and seconds, not samples.
    \item \textbf{Interpretable:} the learned frequencies reveal which
          spectral bands the model finds useful.
    \item \textbf{Fixed token rate:} output tokens are determined by
          \texttt{target\_token\_rate} $\times$ duration, independent of
          input sampling rate.
\end{itemize}

See Appendix~\ref{app:cwt} for a detailed technical description.

\subsubsection{ResampleCNN}
A 1D convolutional stack applied after differentiable resampling to a
fixed token rate. Optionally applies an FFT-based Butterworth anti-aliasing
filter before downsampling.

Key properties:
\begin{itemize}[nosep]
    \item Also handles variable sampling rates via
          \texttt{grid\_sample} resampling.
    \item Learns temporal features directly from the waveform (no
          explicit frequency decomposition).
    \item Comparable parameter count to the CWT approach.
\end{itemize}

See Appendix~\ref{app:cnn} for a detailed technical description.

\subsubsection{Per-Timepoint Linear}
A simple \texttt{nn.Linear} applied independently at each timepoint.
This is the baseline approach. It produces one token per timepoint with
no temporal context.

\subsubsection{Identity}
A passthrough (no temporal transformation). Only meaningful with
spatial session, where the spatial projection already maps channels to
$d_\text{embed}$ dimensions.


%% ===================================================================
\section{Experiments}
\label{sec:experiments}
%% ===================================================================

Experiments are presented in roughly chronological order, grouped by the
question each set of runs was designed to address. Hyperparameters differ
across groups (see Table~\ref{tab:hp-summary}), so absolute performance
numbers should only be compared within a group unless otherwise noted.

\begin{table}[H]
    \centering
    \footnotesize
    \begin{tabular}{llrrr}
        \toprule
        \textbf{Group} & \textbf{Task} & \textbf{LR} & \textbf{WD} & \textbf{Configs $\times$ Folds} \\
        \midrule
        AJILE12\_TOKENIZER\_ABLATION      & Behavior & 6e-4 & 0.09 & 11 $\times$ 2 \\
        AJILE12\_TOKENIZER\_ABLATION\_POSE & Pose     & 6e-4 & 0.09 & 6 $\times$ 2 \\
        AJILE\_NEW\_ALLTOKS               & Behavior & 6e-5 & 0.01 & 11 $\times$ 2 \\
        CWT\_VS\_CNN\_BEHAVIOR            & Behavior & 3e-4 & 0.007 & 4 $\times$ 2 \\
        CWT\_VS\_CNN\_POSE                & Pose     & 3e-4 & 0.007 & 2 $\times$ 2 \\
        \bottomrule
    \end{tabular}
    \caption{Hyperparameter summary per experiment group. Absolute
    performance is comparable only within a group.}
    \label{tab:hp-summary}
\end{table}


% -------------------------------------------------------------------
\subsection{Full Tokenizer Ablation --- Behavior Classification}
\label{sec:ablation-behavior}
% -------------------------------------------------------------------

\textbf{Question:} Across the full design space of channel strategies
and temporal embeddings, which combinations perform best for behavior
classification?

\textbf{Setup:} 11 tokenizer configurations $\times$ 2 folds, using
hyperparameters originally tuned for the per-timepoint linear baseline
(LR=6e-4, WD=0.09). This is a known limitation: these hyperparameters
may not be optimal for all tokenizer families. Token rate was set to
200\,Hz for CWT and ResampleCNN in this group.

\begin{figure}[H]
    \centering
    \includegraphics[width=\textwidth]{ajile12_tokenizer_ablation_bar.pdf}
    \caption{Full tokenizer ablation --- behavior AUROC (mean $\pm$ std
    over 2 folds). Per-channel configurations occupy the right side;
    spatial session configurations cluster on the left.}
    \label{fig:ablation-behavior}
\end{figure}

\textbf{Observations:}
\begin{itemize}[nosep]
    \item Per-channel CWT (dim 512) achieves the highest AUROC
          ($\sim$0.91), followed by per-channel CWT (dim 256) and
          per-channel CNN.
    \item Per-channel configurations consistently outperform their
          spatial session counterparts, suggesting that preserving
          per-electrode identity is valuable for this task.
    \item The per-timepoint linear baseline at dim 256 shows high
          variance across folds and is among the weakest performers,
          confirming the intuition that ignoring temporal structure
          is costly.
    \item Spatial session variants cluster in the 0.84--0.87 range,
          with the ``common layer'' CWT variant performing best among
          them.
\end{itemize}

\textbf{Caveats:} Hyperparameters were tuned for the baseline, which
may disadvantage other architectures. The ResampleCNN used in this
group may not have had anti-aliasing enabled (this was added during the
course of these experiments). Parameter counts were not perfectly
matched between CWT and CNN.


% -------------------------------------------------------------------
\subsection{Full Tokenizer Ablation --- Pose Estimation}
\label{sec:ablation-pose}
% -------------------------------------------------------------------

\textbf{Question:} Do the same patterns hold for a regression task
(pose estimation)?

\textbf{Setup:} 6 tokenizer configurations $\times$ 2 folds, same
hyperparameters as the behavior ablation (LR=6e-4, WD=0.09).

\begin{figure}[H]
    \centering
    \includegraphics[width=\textwidth]{ajile12_tokenizer_ablation_pose_bar.pdf}
    \caption{Full tokenizer ablation --- pose $R^2$ (mean $\pm$ std
    over 2 folds). High variance reflects sensitivity to fold.}
    \label{fig:ablation-pose}
\end{figure}

\textbf{Observations:}
\begin{itemize}[nosep]
    \item Fold variance is very high for pose estimation (e.g.,
          per-channel CWT: fold 0 $R^2 \approx 0.31$, fold 1
          $R^2 \approx 0.07$). This makes it difficult to draw strong
          conclusions from these results.
    \item Spatial CWT (common) achieves the highest mean $R^2$ (0.20),
          but with large standard deviation.
    \item Per-channel CWT and CNN perform similarly, as in the behavior
          task.
    \item Per-timepoint linear is again the weakest performer.
\end{itemize}

\textbf{Caveats:} The high fold variance suggests that pose estimation
results should be treated as preliminary. The fold-1 runs may have
suffered from optimization issues unrelated to the tokenizer.


% -------------------------------------------------------------------
\subsection{Channel Embedding Fusion: Concat vs.\ Add}
\label{sec:concat-vs-add}
% -------------------------------------------------------------------

\textbf{Question:} For per-channel tokenizers, does concatenating the
channel embedding with temporal features outperform additive fusion?

With concatenation, the temporal embedding produces features of dimension
$d_\text{embed} - d_\text{ch}$ (typically $256 - 64 = 192$), and the
channel embedding fills the remaining dimensions. With addition, both
the temporal and channel embeddings are full-dimensional ($d_\text{embed}
= 256$) and are summed. Concatenation was expected to perform better
because it preserves channel identity in a dedicated subspace rather
than mixing it with temporal features.

\textbf{Setup:} Comparison from the AJILE\_NEW\_ALLTOKS group (LR=6e-5,
WD=0.01), which includes both concat and add variants for CWT and CNN.

\begin{figure}[H]
    \centering
    \includegraphics[width=0.65\textwidth]{concat_vs_add_comparison.pdf}
    \caption{Concat vs.\ add fusion for per-channel CWT and CNN. Concat
    is consistently better, but the margin is small.}
    \label{fig:concat-vs-add}
\end{figure}

\textbf{Observations:}
\begin{itemize}[nosep]
    \item Concatenation outperforms addition for both CWT and CNN, but
          the margin is modest ($\sim$0.01--0.02 AUROC).
    \item The advantage of concat may stem from preventing the channel
          identity signal from being ``washed out'' by temporal features
          during summation.
    \item Based on these results, concat fusion was used as the default
          for subsequent per-channel experiments.
\end{itemize}


% -------------------------------------------------------------------
\subsection{Hyperparameter Sensitivity and Sweeps}
\label{sec:hp-sweeps}
% -------------------------------------------------------------------

\textbf{Question:} Are the performance differences observed in the
ablation robust to hyperparameter choice, or are they confounded by
using baseline-tuned hyperparameters for all tokenizers?

A concern with the ablation results is that the learning rate and
weight decay (originally tuned for per-timepoint linear) may happen to
suit some tokenizer families better than others. To address this, we ran
Bayesian hyperparameter sweeps via Weights \& Biases for each major
tokenizer family independently, primarily sweeping learning rate and
weight decay.

\textbf{Key findings:}
\begin{itemize}[nosep]
    \item The optimal learning rate did not vary dramatically across
          tokenizer families. All families performed well in the
          $2\text{e-4}$ to $4\text{e-4}$ range.
    \item This gives some confidence that the ablation results are not
          purely an artifact of hyperparameter mismatch, though it does
          not rule it out entirely.
    \item A shared setting of LR=3e-4, WD=0.007 was identified as a
          good common ground, and was used for the subsequent clean
          head-to-head comparison.
\end{itemize}


% -------------------------------------------------------------------
\subsection{Extended Sweep with Lower Learning Rate}
\label{sec:new-alltoks}
% -------------------------------------------------------------------

\textbf{Question:} With a lower learning rate, do the relative rankings
between tokenizers change?

\textbf{Setup:} 11 configurations (including add-fusion and MLP spatial
variants) $\times$ 2 folds, LR=6e-5, WD=0.01.

\begin{figure}[H]
    \centering
    \includegraphics[width=\textwidth]{ajile_new_alltoks_bar.pdf}
    \caption{Extended tokenizer sweep with lower learning rate ---
    behavior AUROC (mean $\pm$ std over folds).}
    \label{fig:new-alltoks}
\end{figure}

\textbf{Observations:}
\begin{itemize}[nosep]
    \item At this lower learning rate, performance differences between
          tokenizers compress significantly. Most variants cluster in
          the 0.84--0.86 range.
    \item Per-channel CWT and CNN remain at the top, but the gap to
          spatial session variants shrinks.
    \item Per-timepoint linear remains the clear underperformer
          ($\sim$0.79), reinforcing that temporal context is important
          regardless of learning rate.
    \item The MLP spatial projector does not consistently outperform the
          simpler linear projector.
\end{itemize}

\textbf{Interpretation:} The lower learning rate likely prevents all
models from fully converging, compressing performance toward a lower
plateau. The fact that relative rankings are preserved (per-channel
$>$ spatial, CWT $\geq$ CNN $>$ linear) despite the compression is
encouraging for the robustness of these conclusions.


% -------------------------------------------------------------------
\subsection{Clean Head-to-Head: CWT vs.\ ResampleCNN}
\label{sec:cwt-vs-cnn}
% -------------------------------------------------------------------

\textbf{Question:} With matched, well-tuned hyperparameters, how does
the learnable CWT compare to ResampleCNN?

This is the most controlled comparison in this journal. Following the
hyperparameter sweeps, both tokenizers are evaluated with the same
LR (3e-4), WD (0.007), and two embedding dimensions (256 and 512),
using the per-channel strategy with concat fusion.

A deliberate choice was made to reduce the target token rate from
200\,Hz (used in earlier ablations) to 100\,Hz for this group, halving
the token count and corresponding compute. This may account for some
of the reduction in absolute performance compared to the earlier
ablation.

\subsubsection{Behavior Classification}

\begin{figure}[H]
    \centering
    \includegraphics[width=0.85\textwidth]{cwt_vs_cnn_behavior_bar.pdf}
    \caption{CWT vs.\ CNN behavior classification --- best AUROC per
    tokenizer (mean $\pm$ std over folds).}
    \label{fig:cwt-cnn-behavior}
\end{figure}

\begin{figure}[H]
    \centering
    \includegraphics[width=0.85\textwidth]{cwt_vs_cnn_behavior_curves.pdf}
    \caption{Training curves (validation AUROC vs.\ epoch) for
    CWT vs.\ CNN behavior runs (fold 0).}
    \label{fig:cwt-cnn-behavior-curves}
\end{figure}

\subsubsection{Pose Estimation}

\begin{figure}[H]
    \centering
    \includegraphics[width=0.85\textwidth]{cwt_vs_cnn_pose_bar.pdf}
    \caption{CWT vs.\ CNN pose estimation --- best $R^2$ per tokenizer
    (mean $\pm$ std over folds).}
    \label{fig:cwt-cnn-pose}
\end{figure}

\subsubsection{Embedding Dimension Scaling}

\begin{figure}[H]
    \centering
    \includegraphics[width=0.95\textwidth]{cwt_vs_cnn_scaling.pdf}
    \caption{Effect of embedding dimension (256 vs.\ 512) on CWT
    and CNN tokenizers across both tasks.}
    \label{fig:cwt-cnn-scaling}
\end{figure}

\textbf{Observations:}
\begin{itemize}[nosep]
    \item For behavior classification, CWT and CNN perform similarly at
          dim 256 ($\sim$0.88 AUROC), with CWT holding a slight edge.
    \item For pose estimation, CNN takes a slight edge over CWT, though
          both achieve similar $R^2$ values.
    \item At dim 512, both tokenizers show reduced performance compared
          to dim 256 in this setup, possibly due to overfitting at the
          larger embedding size or the interaction with the reduced
          token rate.
    \item Training curves suggest CWT may converge slightly faster in
          early epochs, though both reach similar final performance.
\end{itemize}

\textbf{Interpretation:} The similar performance of CWT and CNN
suggests that the CNN is likely learning some combination of the
frequency-domain features that CWT extracts explicitly. This makes
CWT attractive not because of a clear performance advantage, but
because of its other properties: sampling-rate invariance by
construction, interpretable learned frequencies, and the potential
for regularization effects discussed in Section~\ref{sec:discussion}.


%% ===================================================================
\section{Discussion and Ongoing Conclusions}
\label{sec:discussion}
%% ===================================================================

This section summarizes the tentative conclusions from the experiments
conducted so far. These are working hypotheses, not final claims---the
investigation is ongoing, and several important questions remain open.

\subsection{Per-Channel vs.\ Spatial Session}

Across all experiment groups, per-channel tokenization consistently
outperforms spatial session for behavior classification. The performance
gap is moderate (2--5\% AUROC depending on the group) but consistent.

\textbf{Possible explanations:}
\begin{itemize}[nosep]
    \item Per-electrode identity may carry task-relevant information
          that the spatial projection discards.
    \item The spatial projection is learned per session, which may not
          generalize well with only 2 folds of data.
    \item The Perceiver backbone may benefit from having more tokens
          (per-channel produces $C \times N_\text{time}$) even if they
          are correlated.
\end{itemize}

\textbf{Caveats:} The spatial session approach was always tested with
hyperparameters tuned for the baseline. While the sweeps suggest
optimal learning rates are similar, we cannot rule out that the spatial
approach would close the gap with dedicated architecture tuning
(e.g., different $N_\text{sources}$, projection rank, or training
schedule).

\subsection{CWT vs.\ ResampleCNN}

CWT and ResampleCNN achieve comparable downstream performance across
both tasks. CWT has a slight edge on behavior classification; CNN has
a slight edge on pose estimation. Neither advantage is large enough to
be considered conclusive with the current number of runs.

\textbf{Working hypothesis:} The CNN likely learns to approximate the
same frequency-domain features that CWT extracts explicitly. The value
of CWT may therefore lie not in raw task performance, but in:
\begin{itemize}[nosep]
    \item \textbf{Sampling-rate invariance:} CWT is invariant by
          construction (wavelets are defined in physical units). The CNN
          requires an explicit resampling step to handle variable
          sampling rates.
    \item \textbf{Interpretability:} CWT's learned center frequencies
          reveal which spectral bands the model finds useful, providing
          insight into the data.
    \item \textbf{Potential regularization:} When training across
          datasets of varying quality (different units, noise levels,
          artifacts), the CWT's frequency decomposition may act as a
          learned filter that regularizes the input, making the loss
          landscape more tractable. This hypothesis is motivated by
          preliminary observations from cross-dataset training but
          has not yet been systematically tested.
\end{itemize}

\subsection{Per-Timepoint Linear is Suboptimal}

The per-timepoint linear baseline consistently underperforms all other
temporal embeddings, often by a large margin. This confirms the
intuition that ignoring temporal structure is costly and validates the
effort to explore alternatives.

\subsection{Compute Efficiency}

Reducing the token rate from 200\,Hz to 100\,Hz roughly halves the
number of tokens and corresponding compute. The reduction in task
performance appears modest, suggesting that 100\,Hz may be a
reasonable operating point for many applications. Further investigation
of the performance--compute trade-off at different token rates is
warranted.

\subsection{CWT + CNN Hybrid}

Preliminary experiments with a CWT-compressor architecture (CWT
followed by strided convolutions to further reduce token count) did
not improve over either CWT or CNN alone. The reduced token count may
limit the Perceiver backbone's capacity to learn temporal features.
However, using the CWT as an anti-aliasing or feature-extraction step
\emph{before} a CNN on the resampled features remains an interesting
direction that has not been fully explored.

\subsection{Open Questions}

\begin{itemize}[nosep]
    \item \textbf{Cross-dataset generalization:} All experiments so far
          use a single dataset (AJILE12). The tokenizer comparison may
          yield different results on scalp EEG datasets with different
          noise characteristics and frequency content.
    \item \textbf{Self-supervised pretraining:} The CWT's potential
          regularization properties may be most valuable in a
          self-supervised setting across diverse datasets of varying
          quality. This is a priority for future experiments.
    \item \textbf{Spatial session with more data:} The spatial
          projection may benefit from more training data (more sessions,
          more subjects) to learn a robust channel mixing.
    \item \textbf{Token rate optimization:} A systematic sweep of
          \texttt{target\_token\_rate} would clarify the
          performance--compute trade-off.
    \item \textbf{Learned CWT frequencies:} Analyzing how the CWT's
          center frequencies evolve during training could provide
          insight into which spectral bands are most informative for
          different tasks.
\end{itemize}


%% ===================================================================
\section{Notes \& Observations}
\label{sec:notes}
%% ===================================================================

\begin{itemize}
    \item \textbf{Hyperparameter differences across groups:} The
          experiment groups use different LR/WD settings (see
          Table~\ref{tab:hp-summary}), so absolute performance numbers
          are not directly comparable across groups. Within each group,
          comparisons are fair.

    \item \textbf{Anti-aliasing for ResampleCNN:} The Butterworth
          anti-aliasing filter was added during the course of these
          experiments. Earlier runs (AJILE12\_TOKENIZER\_ABLATION) may
          not have had it enabled, which could have slightly
          disadvantaged the CNN in those comparisons.

    \item \textbf{Parameter count matching:} CWT and CNN were not
          perfectly parameter-matched in all groups. The CWT has
          learnable frequency and cycle-count parameters in addition
          to the projection layer, while the CNN has convolutional
          weights. This is a minor confound but worth noting.

    \item \textbf{Token rate change:} The CWT\_VS\_CNN group uses
          \texttt{target\_token\_rate}=100\,Hz, while the earlier
          ablation used 200\,Hz. This was a deliberate choice to
          reduce compute, but it may account for some of the
          reduction in absolute performance.
\end{itemize}


%% ===================================================================
\begin{appendices}
%% ===================================================================

\section{Learnable CWT Layer}
\label{app:cwt}

This appendix provides a self-contained technical description of the
\texttt{ContinuousCWTLayer} and \texttt{CWTEmbedding} modules. The
layer applies complex Morlet wavelets at learnable center frequencies
to decompose a 1-D signal into a 2-D time--frequency map, then
resamples the time axis to a fixed token count.

\subsection{Learnable Parameters}

The layer has two sets of learnable parameters: \textbf{center
frequencies} $\{f_k\}$ and \textbf{cycle counts} $\{n_k\}$ (one per
frequency bin). Both must remain positive, so the layer stores
unconstrained values and maps them through softplus with a minimum
floor:

\begin{equation}
    f_k = f_\text{min} + \operatorname{softplus}(\tilde{f}_k), \qquad
    n_k = n_\text{min} + \operatorname{softplus}(\tilde{n}_k)
\end{equation}

where $\tilde{f}_k$ and $\tilde{n}_k$ are the unconstrained parameters.
At initialization, the inverse softplus is applied so that the
constrained values start exactly at the user-specified initial values.

\subsection{Morlet Wavelet Construction}

A complex Morlet wavelet at frequency $f$ is a sinusoid windowed by a
Gaussian envelope:

\begin{equation}
    \psi(t) = \exp\!\left(-\frac{t^2}{2\sigma^2}\right) \cdot
    e^{\,i\,2\pi f\, t}, \qquad
    \sigma = \frac{n_\text{cycles}}{2\pi f}
\end{equation}

The Gaussian width $\sigma$ is tied to both the frequency and the
\emph{time--frequency trade-off} parameter $n_\text{cycles}$: more
cycles yield sharper frequency resolution but blurrier time resolution
(Heisenberg uncertainty principle for signal processing).

\begin{figure}[H]
    \centering
    \includegraphics[width=0.85\textwidth]{appendix_cwt_wavelets.pdf}
    \caption{Complex Morlet wavelets at five center frequencies.
    Higher-frequency wavelets are narrower ($\sigma \propto 1/f$),
    while lower-frequency wavelets are wider and contain slower
    oscillations.}
    \label{fig:app-cwt-wavelets}
\end{figure}

\subsection{FFT-Based Convolution}

The wavelet is slid across the signal via cross-correlation, computed
efficiently in the frequency domain:

\begin{enumerate}[nosep]
    \item FFT the signal and the (flipped) wavelet kernel.
    \item Pointwise multiply in the frequency domain.
    \item Inverse FFT and extract the valid region.
\end{enumerate}

Because the wavelet kernels depend on the per-batch-item sampling rate
$f_s$ (via $t_\text{sec} = t_\text{samples} / f_s$), different
sampling rates within a batch are handled naturally without
approximation.

\begin{figure}[H]
    \centering
    \includegraphics[width=0.85\textwidth]{appendix_cwt_scalogram.pdf}
    \caption{CWT magnitude (scalogram) of a synthetic chirp signal
    (2--20\,Hz sweep) with a 10\,Hz burst at $t = 1.5$--$2.0$\,s.
    Each wavelet picks up energy where its frequency is present.}
    \label{fig:app-cwt-scalogram}
\end{figure}

\subsection{Time Resampling}

Signals in a batch can have different lengths (zero-padded to
$\text{Max\_T}$). The time--frequency map is treated as a 2-D image
(height = frequency bins, width = time) and resampled to a fixed
\texttt{target\_time\_tokens} count using \texttt{grid\_sample} with
bilinear interpolation. Per-item sampling coordinates are built so
that only the valid (non-padded) region is covered.

\begin{figure}[H]
    \centering
    \includegraphics[width=0.9\textwidth]{appendix_cwt_resampling.pdf}
    \caption{Variable-length signals resampled to the same number of
    output tokens. The full 3\,s signal and a 1.5\,s signal both
    produce 64 tokens, but each covers only its valid duration.}
    \label{fig:app-cwt-resampling}
\end{figure}

\subsection{Magnitude and Phase Extraction}

The complex CWT output is converted to polar form:

\begin{align}
    \text{magnitude} &= \sqrt{\text{re}^2 + \text{im}^2 + \epsilon} \\
    \text{phase} &= \frac{\operatorname{atan2}(\text{im}, \text{re})}{\pi}
    \cdot w, \qquad
    w = \frac{|\mathbf{z}|^2}{|\mathbf{z}|^2 + \epsilon_\text{phase}}
\end{align}

Phase is suppressed for near-zero magnitudes via the soft weight $w$,
preventing noisy gradients where no frequency content is present.

\begin{figure}[H]
    \centering
    \includegraphics[width=0.65\textwidth]{appendix_cwt_phase_stability.pdf}
    \caption{Phase stability weighting. Where magnitude is large,
    $w \approx 1$ and phase passes through unchanged. Near zero
    magnitude, phase is smoothly driven toward zero.}
    \label{fig:app-cwt-phase-stability}
\end{figure}

\subsection{Sampling-Rate Invariance}

A key property: the same physical signal sampled at different rates
produces approximately the same output, because wavelets are defined
in physical units (Hz, seconds) and the resampling maps valid duration
to a fixed token count.

\begin{figure}[H]
    \centering
    \includegraphics[width=0.95\textwidth]{appendix_cwt_fs_invariance.pdf}
    \caption{The same 3\,s chirp signal sampled at 128, 256, and
    512\,Hz produces near-identical scalograms. Differences are due
    to discretization only.}
    \label{fig:app-cwt-fs-invariance}
\end{figure}

\subsection{Frequency Selectivity}

Each wavelet acts as a bandpass filter centered at its learnable
frequency. The selectivity (bandwidth) is controlled by
$n_\text{cycles}$: more cycles yield a narrower, more selective
response.

\begin{figure}[H]
    \centering
    \includegraphics[width=0.85\textwidth]{appendix_cwt_freq_response.pdf}
    \caption{Frequency response of each wavelet. Pure sine waves at
    known frequencies are passed through the layer; the response
    confirms each wavelet is a bandpass filter peaked at its center
    frequency.}
    \label{fig:app-cwt-freq-response}
\end{figure}

\begin{figure}[H]
    \centering
    \includegraphics[width=0.85\textwidth]{appendix_cwt_ncycles.pdf}
    \caption{Effect of $n_\text{cycles}$ on frequency selectivity
    (10\,Hz wavelet shown). More cycles yield a narrower bandwidth.}
    \label{fig:app-cwt-ncycles}
\end{figure}

\subsection{Output Shape}

The full \texttt{CWTEmbedding} pipeline:
\begin{enumerate}[nosep]
    \item \texttt{ContinuousCWTLayer}: $(B, S, T) \to (B, S, 2, F, T')$
          where $S$ = sources, $F$ = frequency bins, $T'$ = target time
          tokens, and the 2 channels are magnitude and phase.
    \item Flatten: $(B, T', S \cdot 2 \cdot F)$
    \item Linear projection: $(B, T', d_\text{embed})$
\end{enumerate}


%% -------------------------------------------------------------------
\section{ResampleCNN Architecture}
\label{app:cnn}
%% -------------------------------------------------------------------

This appendix describes the \texttt{ResampleCNNEmbedding} module, which
serves as the primary non-CWT temporal embedding.

\subsection{Pipeline Overview}

\begin{figure}[H]
    \centering
    \includegraphics[width=0.9\textwidth]{appendix_cnn_pipeline.pdf}
    \caption{ResampleCNN processing pipeline. The signal is optionally
    low-pass filtered, resampled to a fixed token count, processed by
    a convolutional stack, and linearly projected to the embedding
    dimension.}
    \label{fig:app-cnn-pipeline}
\end{figure}

\subsection{Anti-Aliasing Filter}

Before downsampling, an FFT-based Butterworth low-pass filter removes
frequencies above the target Nyquist to prevent aliasing. The filter is
applied in the frequency domain:

\begin{equation}
    |H(f)| = \frac{1}{\sqrt{1 + (f / f_c)^{2n}}}
\end{equation}

where $f_c$ is the per-item cutoff frequency (proportional to the
downsampling ratio) and $n = 6$ is the filter order. The filter is a
no-op when the signal is being upsampled.

\begin{figure}[H]
    \centering
    \includegraphics[width=0.85\textwidth]{appendix_cnn_butterworth.pdf}
    \caption{Butterworth low-pass filter magnitude response at different
    downsampling ratios. Higher ratios require more aggressive
    filtering.}
    \label{fig:app-cnn-butterworth}
\end{figure}

\subsection{Differentiable Resampling}

The signal is resampled to a fixed number of time tokens using
\texttt{grid\_sample}. The 1-D signal is treated as a single-row image,
and bilinear interpolation along the width (time) axis provides smooth,
differentiable resampling. Per-item grid coordinates map only the valid
(non-padded) portion of each signal.

\begin{figure}[H]
    \centering
    \includegraphics[width=0.85\textwidth]{appendix_cnn_resampling.pdf}
    \caption{Anti-aliased resampling of a chirp signal (2--60\,Hz) to
    different token counts. The Butterworth filter prevents aliasing
    artifacts in the downsampled versions.}
    \label{fig:app-cnn-resampling}
\end{figure}

\subsection{Convolutional Stack and Projection}

After resampling, a stack of 1-D convolutions (default: 2 layers,
kernel size 9, 64 filters, GELU activation, same-padding) extracts
temporal features. A final linear layer projects from the filter
dimension to the model's embedding dimension. The output shape is
$(B, T', d_\text{embed})$, matching the CWT embedding's interface.

\end{appendices}

\end{document}
